\section{Thermodynamic limits at fixed spacing}
\label{sec:thermo}

All parameters are fixed in this section. We distinguish
unrestricted open-graph minima, physical zero-exterior-flux
open minima, and periodic physical minima.

\subsection{Open graphs and a bounded cutting estimate}

For a rectangular $Y\subset\mathbb Z^3$, keep internal links and
plaquettes. The open fermion matrix is $m\beta+\sum_{e\subset Y}h_e$,
where for $e=(x,y=x+e_j)$
\begin{equation}
 h_e=\frac{\alpha_j}{2ia}
 (\overline U_e|x\rangle\langle y|-U_e|y\rangle\langle x|)
 +\frac{r\beta}{2a}
 (|x\rangle-U_e|y\rangle)(\langle x|-\overline U_e\langle y|).
 \label{eq:open-edge}
\end{equation}
The Wilson diagonal at $x$ is $m+r\deg_Y(x)/(2a)$,
not the periodic value at boundary sites. Its mass coefficient
is still at least $m$. Denote the auxiliary minimum by
$e_{\rm op}(Y)$ and the physical minimum with
$\diver_YE+q=\eta$ by $E_Y^{\rm op}(\eta)$.
There are no exterior link rotors: this convention fixes exterior
flux to zero.

The proof of Theorem~\ref{thm:main} applies on this finite graph:
its incidence matrix replaces the periodic one, the Wilson mass
is positive, and the internal rotor kernels are unchanged.
For connected $Y$ the admissibility condition is
$|\sum_Y\eta|\le2|Y|$, and
\begin{equation}
 e_{\rm op}(Y)=E_Y^{\rm op}(0),\qquad
 E_Y^{\rm op}(\eta)\ge e_{\rm op}(Y)+m|\sum_Y\eta|.
 \label{eq:open-comparison}
\end{equation}
The empty-link case is the same finite-CAR argument.
Uniformly for rectangles of degree at most six,
\begin{equation}
 -2h_*|Y|\le e_{\rm op}(Y)\le3|Y|/(g^2a).
 \label{eq:open-extensive}
\end{equation}
The upper trial has zero electric flux and two negative-$\beta$
spinors at every site; discard its nonpositive fermion energy.

\begin{lemma}[Cutting with retained rotors]
\label{lem:cut}
Select rectangular open blocks partitioning the sites of a cubic
graph. Let $\mathcal C$ be all original links not internal to
these open blocks, including periodic closing links even when
both endpoints lie in the same block, and put $n_c=|\mathcal C|$.
On the original auxiliary space define
\[
 H_{\rm cut}=\sum_B H_B^{\rm op}
                       +c_E\sum_{e\in\mathcal C}E_e^2,
 \qquad c_E=g^2/(2a).
\]
Then
\begin{equation}
 \|H-H_{\rm cut}\|\le Bn_c,\qquad
 B=(4r+2+8/g^2)/a,
 \label{eq:cut-bound}
\end{equation}
and the cut auxiliary minimum equals $\sum_B e_{\rm op}(B)$.
The same statement holds when adjoining periodic closing links
to an open cube, after adding their free rotor factors.
\end{lemma}
\begin{proof}
The endpoint matrix \eqref{eq:open-edge} has norm at most
$r/a+1/(2a)$ and anticommutes with $C_5$. Its eight-dimensional
spectrum is paired, so
$\|d\Gamma(h_e)\|=\frac12\Tr|h_e|\le(4r+2)/a$.
A plaquette multiplier has norm at most $2/(g^2a)$ and a link
meets at most four plaquettes. Remove the cut fermion terms and
every plaquette meeting a cut link; counting a plaquette more
than once only enlarges the estimate.
The unbounded electric term on each cut link is retained as an
independent $L^2(U(1))$ rotor with minimum zero.
Even fermion operators on disjoint blocks commute under the graded
Fock decomposition, so their unrestricted minima add.
The variational principle gives the claimed spectral comparison.
\end{proof}

\subsection{The vacuum limit}

Write $Q_n=\{0,\ldots,n-1\}^3$, $e_{\rm op}(n)=e_{\rm op}(Q_n)$,
and let $E_0^{\rm per}(n)$ be the physical periodic vacuum energy.

\begin{theorem}[Vacuum thermodynamic limit]
\label{thm:vacuum-limit}
The finite limit
\begin{equation}
 e_{\vac}=\lim_{n\to\infty}\frac{e_{\rm op}(n)}{n^3}
         =\lim_{n\to\infty}\frac{E_0^{\rm per}(n)}{n^3}
 \label{eq:vacuum-limit}
\end{equation}
exists, and for every integer $\ell\ge3$,
\begin{equation}
 \left|e_{\vac}-\frac{e_{\rm op}(\ell)}{\ell^3}\right|
 \le\frac{3B}{\ell},\qquad
 \left|e_{\vac}-\frac{E_0^{\rm per}(\ell)}{\ell^3}\right|
 \le\frac{6B}{\ell}.
 \label{eq:vacuum-rate}
\end{equation}
Physical energy per spatial volume is $e_{\vac}/a^3$.
\end{theorem}
\begin{proof}
For $n=k\ell$, cutting an open cube into $k^3$ cubes removes
$3(k-1)n^2\le3n^3/\ell$ links. Hence
$|e_{\rm op}(n)-k^3e_{\rm op}(\ell)|\le3Bn^3/\ell$.
For arbitrary $n$, subdivide each coordinate into intervals of
length $\ell$ and one possible remainder. Remainder boxes contain
at most $3\ell n^2$ sites; their energies obey
\eqref{eq:open-extensive}. The cut-link count is at most
$3n^3/\ell$. Thus at every fixed $\ell$,
\[
 \frac{e_{\rm op}(\ell)}{\ell^3}-\frac{3B}{\ell}
 \le\liminf_n\frac{e_{\rm op}(n)}{n^3}
 \le\limsup_n\frac{e_{\rm op}(n)}{n^3}
 \le\frac{e_{\rm op}(\ell)}{\ell^3}+\frac{3B}{\ell}.
\]
The limsup--liminf difference is at most $6B/\ell$ for every
$\ell$, proving existence and the first estimate.
Closing an open $n$-cube adds $3n^2$ links. Lemma~\ref{lem:cut}
gives $|E_{\rm aux}^{\rm per}(n)-e_{\rm op}(n)|\le3Bn^2$.
Theorem~\ref{thm:main} identifies
$E_{\rm aux}^{\rm per}(n)=E_0^{\rm per}(n)$, completing the proof.
\end{proof}

This cutting proof concerns the auxiliary tensor product.
It does not set the boundary flux of an arbitrary physical
state to zero.

\subsection{A fixed periodic source cell}

Choose a cubic period cell of side $p\ge1$, with $V_c=p^3$ sites,
$K_c\ge1$ labeled positive nuclei and charge $Z_c\le2V_c$.
A rectangular period can be replaced by a cubic supercell.
All periodic source sizes satisfy $n\in p\mathbb N$ and $n\ge3$.
Choose integers $q_x^{(c)}\in\{0,1,2\}$ with total $Z_c$
and an internal integer flow
\begin{equation}
 \diver F=\nu-q^{(c)},\qquad
 \mathcal F_c=\sum_{e\ {\rm in\ cell}}F_e^2.
 \label{eq:cell-flow}
\end{equation}
A spanning tree provides it; a one-site cell has zero right side.
Set $\mu=m+3r/a$ and $\rho=Z_c/V_c$.

\begin{proposition}[Fixed-cell relative bounds]
\label{prop:cell}
For $N_c=(n/p)^3$, $K=K_cN_c$ and $Z_{\rm tot}=Z_cN_c$,
\begin{equation}
 0\le E_\nu^{\rm per}(n)-E_0^{\rm per}(n)-mZ_{\rm tot}
 \le C_{\rm cell}K,
 \label{eq:cell-bound}
\end{equation}
where
\begin{equation}
 C_{\rm cell}
 =\frac{3rZ_c+(6+6r+3/g^2)V_c+(g^2/2)\mathcal F_c}{aK_c}.
 \label{eq:cell-constant}
\end{equation}
The constant may depend on the cell and its screening flow.
Every periodic source ground state also satisfies
\begin{equation}
 \langle H_E\rangle\le A_c n^3,\qquad
 A_c=2(h_*-\mu)+\mu\rho+c_E\mathcal F_c/V_c+3/(g^2a).
 \label{eq:electric-budget}
\end{equation}
\end{proposition}
\begin{proof}
Repeat the cell flow, put intercell links in electric flux zero,
and occupy the two negative-$\beta$ spinors plus
$q_x^{(c)}$ positive-$\beta$ spinors at each site.
All Gauss constraints hold and hopping expectations vanish.
Its periodic energy is at most
\begin{equation}
 -2\mu n^3+\mu Z_{\rm tot}
                 +c_EN_c\mathcal F_c+3n^3/(g^2a).
 \label{eq:cell-trial}
\end{equation}
Subtract $E_0^{\rm per}\ge-2h_*n^3$ and $mZ_{\rm tot}$
to obtain the upper bound. The lower bound is
\eqref{eq:main}. Finally $H_B\ge0$ and the fermion lower bound
give $\langle H_E\rangle\le E_\nu^{\rm per}+2h_*n^3$;
insert \eqref{eq:cell-trial}.
\end{proof}

The budget \eqref{eq:electric-budget} is absolute and extensive,
not a bound on electric excitations relative to the vacuum.
On an open rectangular union of complete cells, the same trial
gives
\begin{equation}
 E_Y^{\rm op}(\nu)\le
 \left[-2m+m\rho+c_E\mathcal F_c/V_c+3/(g^2a)\right]|Y|.
 \label{eq:open-cell-trial}
\end{equation}
Indeed the site's diagonal expectation is
$(q_x^{(c)}-2)(m+r\deg_Y(x)/(2a))\le(q_x^{(c)}-2)m$.

\begin{theorem}[Open zero-flux source limit]
\label{thm:open-source}
For each fixed cell as above,
\begin{equation}
 e_\nu^{\rm op}
 =\lim_{\substack{n\to\infty\\p\mid n}}
       \frac{E_{Q_n}^{\rm op}(\nu)}{n^3}
 =\inf_{\substack{\ell\in p\mathbb N\\\ell\ge3}}
   \left\{\frac{E_{Q_\ell}^{\rm op}(\nu)}{\ell^3}
                      +\frac{3B}{\ell}\right\}
 \label{eq:open-source-limit}
\end{equation}
exists and is finite. In addition,
\begin{equation}
 e_{\vac}+m\rho
 \le\liminf_{\substack{n\to\infty\\p\mid n}}\frac{E_\nu^{\rm per}(n)}{n^3}
 \le\limsup_{\substack{n\to\infty\\p\mid n}}\frac{E_\nu^{\rm per}(n)}{n^3}
 \le e_\nu^{\rm op}.
 \label{eq:periodic-bracket}
\end{equation}
\end{theorem}
\begin{proof}
Take products of zero-exterior-flux physical block ground states
and give newly joining rotors electric flux zero. This state
satisfies all source Gauss constraints on the joined graph.
Lemma~\ref{lem:cut} gives the upper almost-subadditivity
\[
 E_{Q_{k\ell}}^{\rm op}(\nu)
 \le k^3E_{Q_\ell}^{\rm op}(\nu)+3B(k\ell)^3/\ell.
\]
For general $n,\ell\in p\mathbb N$, the rectangular remainders
are unions of whole admissible cells. Their trial energy is
$O(\ell n^2)$ by \eqref{eq:open-cell-trial}. Consequently
the limsup of $E_{Q_n}^{\rm op}(\nu)/n^3$ is at most
$E_{Q_\ell}^{\rm op}(\nu)/\ell^3+3B/\ell$ for every fixed
compatible $\ell$. Let $\gamma$ be their infimum.
By definition $E_{Q_n}^{\rm op}(\nu)/n^3+3B/n\ge\gamma$;
hence the liminf is at least $\gamma$ and the limsup at most
$\gamma$. The extensive lower bound makes the limit finite.
Only one-sided almost-subadditivity was used.

Adjoining zero-flux closing rotors gives
$E_\nu^{\rm per}(n)\le E_{Q_n}^{\rm op}(\nu)+3Bn^2$.
This proves the last inequality in \eqref{eq:periodic-bracket}.
Its first follows from the charge-cost and vacuum-limit theorems.
\end{proof}

Thus the \emph{open-boundary} relative energy per nucleus has
limit $(V_c/K_c)(e_\nu^{\rm op}-e_{\vac})$.
A fixed subtraction $S_\infty$ subtracts its fixed density.
The periodic subtraction $S_n$ has the same limiting density:
$\G_n(0)\to\G_\infty(0)$ by ordinary Fourier Riemann convergence
away from zero and the uniform $O(\delta)$ shell bound within
$0<|k|\le\delta n$, followed by $\delta\downarrow0$.

\subsection{The periodic-source question}

We state precisely why \eqref{eq:periodic-bracket} is not yet
a periodic-source limit theorem. For a cut set $\mathcal C$,
orient each cut link once, write $f\in\mathbb Z^{\mathcal C}$,
and let $b_B(f)_x$ be its signed outgoing flux at a boundary
site $x$ of block $B$. Internal Gauss law becomes
$\diver_BE+q=\nu|_B-b_B(f)$.
Diagonalizing the retained free cut rotors gives the exact comparison
\begin{equation}
 \left|E_\nu-\inf_f
 \left\{\sum_B E_B^{\rm op}(\nu|_B-b_B(f))
                           +c_E\sum_{e\in\mathcal C}f_e^2\right\}\right|
 \le B|\mathcal C|.
 \label{eq:flux-decomposition}
\end{equation}
An empty block sector has energy $+\infty$.
At fixed $f$ the physical space of the cut Hamiltonian is the
tensor product of the displayed block sectors. The direct sum
over $f$ gives its minimum, and the bounded-cut estimate proves
\eqref{eq:flux-decomposition}. This step retains the constraints;
replacing all block energies by their vacuum minima would lose
the source comparison.

A sufficient additional hypothesis on commensurate cubes is
\begin{equation}
 E_B^{\rm op}(\nu-b_B(f))\ge E_B^{\rm op}(\nu)
      -C_\partial|\partial_{\rm links}B|
      -C_f\sum_{e\in\partial_{\rm links}B}f_e^2,
 \tag{BF}\label{eq:bf}
\end{equation}
with finite constants independent of side and flux, whenever
the left sector is admissible. Boundary links here are counted
as incidences, so a link shared by two blocks occurs twice.
Hypothesis (BF) is \emph{not proved}.

For clarity, its sufficiency can be checked without a restrictive
smallness condition on $C_f$. Average the grid of $\ell$-cubes
over translations by whole period cells. If $\ell\mid n$, a
given link is cut with probability at most $p/\ell$.
By \eqref{eq:electric-budget}, some grid therefore obeys
\[
 \left\langle\sum_{\rm cut}E_e^2\right\rangle
       \le pA_cn^3/(c_E\ell).
\]
There are $3n^3/\ell$ cut links. Applying (BF) in the state's
flux decomposition would give
\begin{equation}
 \frac{E_\nu^{\rm per}(n)}{n^3}\ge
 \frac{E_{Q_\ell}^{\rm op}(\nu)}{\ell^3}
 -\frac{3B+6C_\partial}{\ell}
 -\frac{(2C_f-c_E)_+pA_c}{c_E\ell}.
 \label{eq:conditional}
\end{equation}
For general compatible $n$, use whole-cell rectangular
remainders and average the global origin. Cut probability is
at most $p/\ell+p/n$; the remaining $O(\ell n^2)$ sites have
the absolute extensive lower bound. Thus \eqref{eq:conditional}
persists for the liminf at fixed $\ell$. Sending $\ell\to\infty$
would prove equality of the periodic and open limits.
Any finite $C_f$ suffices. Direct absorption without averaging
would instead require $2C_f\le c_E$, because of double incidence.

Admissibility in such a comparison matters. For example take
$\nu_x=4$ on even $x_1$ layers and zero on odd layers. Its mean
is the admissible value two per site. An odd $\ell$-cube starting
on an even layer has total source charge
$2\ell^3+2\ell^2>2|B|$, so its zero-flux sector is empty.
Yet $q_x=2$, internal $E_{x,1}=(-1)^{x_1}$, other components
zero, and outward flux $+1$ on both $x_1$ faces give a
finite-energy state. No finite boundary cost compares this
state to the empty zero-flux sector. Whole-cell cubes avoid
this obstruction but do not by themselves establish (BF).

The averaging estimate controls exposed electric energy; it
does not repair the block Gauss sectors. The Graf--Schenker
continuum Coulomb localization inequality, as stated in
\cite[Theorem 2]{ThermoII}, likewise does not supply (BF) for
compact integer fluxes. The periodic-source limit remains
conditional on a screening or gluing input of this kind.

\subsection{Limits of the conclusions}

The results are for a fixed spatial regulator with nonnegative
\emph{bare} Wilson mass. They do not establish that a
renormalization trajectory with fixed physical parameters stays
in that range. The self-energy bound diverges as $a^{-1}$.
There is no construction of renormalized continuum QED or
cutoff-uniform stability constant.
Compact gauge rotors provide a different finite regulator from
a transverse photon Fock-space momentum cutoff. No Coulomb-phase
or gapless-photon assertion is needed.

The electric kernel used here is that of the unshifted rotor.
Additional flux projections, nonintegral background electric
shifts, or finite electric truncations need not preserve the
positive measure in the proof and are not covered.
The charge gap compares different total-charge sectors; binding,
neutral pair gaps and crystal equilibrium require further
arguments. Finally, thermodynamic limits for the vacuum and for
open zero-flux sources do not imply the unproved periodic-source
limit. These distinctions are part of the theorem's scope.
